%-------------------------------------------------------------------------------
% Hridik Hingorani: one-page resume (hosted as public/media/Hridik_Hingorani_Resume.pdf)
% Layout follows Jake Gutierrez's resume template (MIT), with navy section headings and
% underlined contact links. Recreated Nov 2026 from Hridik_Hingorani_Resume-4.pdf.
% Build: pdflatex resume.tex
%-------------------------------------------------------------------------------
\documentclass[letterpaper,11pt]{article}

\usepackage{latexsym}
\usepackage[empty]{fullpage}
\usepackage{titlesec}
\usepackage[usenames,dvipsnames]{xcolor}
\usepackage{enumitem}
\usepackage[hidelinks]{hyperref}
\usepackage{fancyhdr}
\usepackage[english]{babel}
\usepackage{tabularx}
\usepackage{amsmath}
\usepackage[T1]{fontenc}
\usepackage{lmodern}
\input{glyphtounicode}

\definecolor{heading}{HTML}{1F3A68}

\pagestyle{fancy}
\fancyhf{}
\fancyfoot{}
\renewcommand{\headrulewidth}{0pt}
\renewcommand{\footrulewidth}{0pt}

\addtolength{\oddsidemargin}{-0.6in}
\addtolength{\evensidemargin}{-0.6in}
\addtolength{\textwidth}{1.2in}
\addtolength{\topmargin}{-0.7in}
\addtolength{\textheight}{1.4in}

\urlstyle{same}
\raggedbottom
\raggedright
\setlength{\tabcolsep}{0in}

\titleformat{\section}{\vspace{-6pt}\color{heading}\bfseries\raggedright\large}{}{0em}{}[\color{heading}\titlerule\vspace{-5pt}]
\titlespacing*{\section}{0pt}{10pt}{9pt}
\pdfgentounicode=1

\newcommand{\ulink}[2]{\href{#1}{\underline{#2}}}
\newcommand{\resumeItem}[1]{\item{\fontsize{9.5pt}{11.2pt}\selectfont #1\par}\vspace{-1pt}}
\newcommand{\resumeSubheading}[4]{
  \vspace{-2pt}\item
    \begin{tabular*}{\textwidth}[t]{l@{\extracolsep{\fill}}r}
      \textbf{#1} & \textbf{#2} \\
      \textit{\small#3} & \textit{\small #4} \\
    \end{tabular*}\vspace{-7pt}
}
\newcommand{\resumeProjectHeading}[2]{
    \vspace{-1pt}\item
    \begin{tabular*}{\textwidth}{l@{\extracolsep{\fill}}r}
      \small#1 & \textbf{\small #2} \\
    \end{tabular*}\vspace{-7pt}
}
\newcommand{\resumeSubHeadingListStart}{\begin{itemize}[leftmargin=0in, label={}]}
\newcommand{\resumeSubHeadingListEnd}{\end{itemize}}
\newcommand{\resumeItemListStart}{\begin{itemize}[leftmargin=0.32in, label=\textbullet, itemsep=0pt, topsep=2pt]}
\newcommand{\resumeItemListEnd}{\end{itemize}\vspace{-5pt}}

\begin{document}

\begin{center}
    {\Huge \scshape Hridik Hingorani} \\ \vspace{3pt}
    \small (447) 902-5994 \enspace$|$\enspace \ulink{mailto:hridikh2@illinois.edu}{hridikh2@illinois.edu} \enspace$|$\enspace
    \ulink{https://linkedin.com/in/HridikH}{linkedin.com/in/HridikH} \enspace$|$\enspace
    \ulink{https://github.com/HridikH}{github.com/HridikH} \enspace$|$\enspace
    \ulink{https://hridikhingorani.com}{hridikhingorani.com}
\end{center}
\vspace{-12pt}

%-----------EDUCATION-----------
\section{Education}
  \resumeSubHeadingListStart
    \resumeSubheading
      {University of Illinois Urbana-Champaign}{Expected May 2027}
      {B.S. in Computer Engineering, Minor in Mathematics}{Urbana, IL}
      \resumeItemListStart
        \resumeItem{\textbf{Coursework:} Control Systems, Robot Dynamics and Control, Introduction to Robotics, Principles of Safe Autonomy, Operating Systems, Deep Learning for Computer Vision, Digital Systems Lab}
      \resumeItemListEnd
  \resumeSubHeadingListEnd

%-----------EXPERIENCE-----------
\section{Experience}
  \resumeSubHeadingListStart
    \resumeSubheading
      {Software Associates}{May 2026 -- Aug. 2026}
      {AI Software Engineering Intern}{Remote}
      \resumeItemListStart
        \resumeItem{Built an algorithmic trading system in Python from scratch, combining five technical indicators (RSI, MACD, Bollinger Bands, moving averages, ATR) into a weighted signal engine that issued automated buy, sell, and hold decisions}
        \resumeItem{Integrated the Interactive Brokers API for live market data and order execution with automated stop-loss, take-profit, and trailing-stop logic, and shipped a Flask dashboard streaming live signals, open positions, and P\&L}
      \resumeItemListEnd
    \resumeSubheading
      {University of Illinois Urbana-Champaign}{Jan. 2026 -- Present}
      {Course Assistant, ECE 385: Digital Systems Laboratory}{Urbana, IL}
      \resumeItemListStart
        \resumeItem{Lead lab sections for 60+ students on SystemVerilog, FPGA timing closure, and hardware-software co-design}
        \resumeItem{Debug failures across the stack, from simulation waveforms to on-board behavior, isolating whether a fault lives in the RTL, the constraints, or the hardware}
      \resumeItemListEnd
    \resumeSubheading
      {HealthArk Insights}{May 2025 -- Aug. 2025}
      {AI Software Development Intern}{Mumbai, India}
      \resumeItemListStart
        \resumeItem{Shipped three production Python features to a client-facing document platform (summarization, web scraping, dynamic formatting), holding 92\% extraction accuracy across the document simplification pipeline}
        \resumeItem{Gathered and validated requirements directly with clients, taking each feature from specification through production release}
      \resumeItemListEnd
  \resumeSubHeadingListEnd

%-----------PROJECTS-----------
\section{Projects}
  \resumeSubHeadingListStart
    \resumeProjectHeading
      {\textbf{Humanoid Push Recovery, Booster K1} $|$ \emph{Python, PPO, mjlab, MuJoCo, Slurm}}{Fall 2026}
      \resumeItemListStart
        \resumeItem{Deployed the policy sim-to-real on K1 hardware at 50 Hz: zero falls across 4 gantry runs and about 18 hand pushes, peak tilt 32 degrees against a fall threshold of about 60 degrees, with the K1's built-in fall safeguard never engaging}
        \resumeItem{Cut falls from 39\% to 13\% on unseen 15--40 N$\cdot$s pushes by replacing a failing curriculum with a survival-gated 13-level push ladder and domain randomization}
        \resumeItem{Trained a PPO walking policy for a 19.7 kg humanoid across 4,096 parallel simulations on NCSA Delta, modeling pushes as ramping torso impulses}
      \resumeItemListEnd
    \resumeProjectHeading
      {\textbf{Self-Driving Stack, Polaris GEM e4} $|$ \emph{C++, Python, ROS 2, Gazebo}}{Fall 2026}
      \resumeItemListStart
        \resumeItem{Owned dynamics, controls, and testing on a four-person team building a ROS 2 autonomy stack for a full-scale EV: lane following, stops for signs, pedestrians, and obstacles, and auto-resume}
        \resumeItem{Wrote the control layer: Pure Pursuit with a Stanley baseline, curvature-based speed planning, a speed PID, and actuator limits, stopping within 2 m of the goal}
        \resumeItem{Built a track simulator and Gazebo world from recorded ROS bags for controller testing}
      \resumeItemListEnd
    \resumeProjectHeading
      {\textbf{Real-Time Task-Space Impedance Control, CRS Arm} $|$ \emph{C, TI C2000 (TMS320F28335)}}{Spring 2026}
      \resumeItemListStart
        \resumeItem{Implemented a 1 kHz task-space impedance controller in bare-metal C, reading joint encoders and sending joint torque commands directly to the motor amplifiers of a 3-joint arm}
        \resumeItem{Added per-joint Coulomb and viscous friction compensation, per-waypoint stiffness modes for selective compliance, and torque saturation; validated on peg-in-hole insertion}
        \resumeItem{Commanded contact force open-loop from estimated contact stiffness with no force sensor, and verified it on a scale}
      \resumeItemListEnd
    \resumeProjectHeading
      {\textbf{RISC-V Preemptive Multitasking OS Kernel} $|$ \emph{C, RISC-V Assembly, QEMU, GDB}}{Spring 2026}
      \resumeItemListStart
        \resumeItem{Root-caused a page-table entry bug in a $\sim$6,000-line kernel by stepping the fault path in GDB against QEMU and inspecting translations by hand}
        \resumeItem{Wrote the preemptive scheduler, Sv39 virtual memory, VirtIO block, UART, and RTC drivers, a FAT filesystem, and fork, exec, and pipes behind a shell}
      \resumeItemListEnd
    \resumeProjectHeading
      {\textbf{UR3 Vision-Guided Contour Drawing} $|$ \emph{Python, ROS, OpenCV}}{Spring 2026}
      \resumeItemListStart
        \resumeItem{Derived closed-form inverse kinematics for the 6-DOF UR3 with product-of-exponentials forward kinematics, feeding a Canny and contour pipeline that turns images into stroke paths}
        \resumeItem{Hit joint targets within 0.0005 rad over ROS and drew a 142-stroke image on hardware in 7 minutes}
      \resumeItemListEnd
  \resumeSubHeadingListEnd

%-----------SKILLS-----------
\section{Technical Skills}
 \begin{itemize}[leftmargin=0in, label={}]
    \fontsize{9.5pt}{11.6pt}\selectfont{\item{
     \textbf{Languages}{: C, C++, Python, MATLAB, SystemVerilog, RISC-V Assembly} \\
     \textbf{Controls}{: PID, impedance control, LQR, pole placement and observers, Pure Pursuit/Stanley, analytic IK, PoE kinematics} \\
     \textbf{Learning \& Simulation}{: PPO, domain randomization, mjlab, MuJoCo, Gazebo, Slurm} \\
     \textbf{Embedded \& Tools}{: TI C2000, ESP32, bare-metal firmware, I2C, UART, ADC, PWM, ROS 2, Linux, Git, GDB, QEMU, KiCad}
    }}
 \end{itemize}

\end{document}
